% Appendix: Labb 1, from labs/lab1/README.md. The guide's headings are sections, and its tasks,
% which are headings of their own in the guide, are collected under the section Uppgifterna.
%
% Deliberate differences: the lecture links are chapter and section references, the common
% safety rules the guide points to in labs/README.md are in the preface, and the table in task 1
% has two empty rows more than the guide's, for a station with more than two contact elements
% per button (the guide says to add rows).
%
% \labprep[title]{F1}{label} is a local helper: dtbook.sty's \prep sets no link target, and the
% preparations are referred to as lab1:f1 to lab1:f4, with \ref printing "F1". It is
% \providecommand, since lab2.tex defines the same one.

\makeatletter
\providecommand{\labprep}[3][]{\prep{#2\if\relax\detokenize{#1}\relax\else\ -- #1\fi}%
  \phantomsection\def\@currentlabel{#2}\label{#3}}
\makeatother

\chapter{Labb 1: Kopplingsnät}
\label{app:lab1}

I den här laborationen bygger du logiska funktioner av strömkontakter: två tryckknappar, en lampa
och ett aggregat på 24~V. Samma funktioner som i kapitel~\ref{c2:ch} byggdes av grindar blir här
serie- och parallellkopplingar av slutande och brytande kontakter, och lampans läge är utgången.

Laborationen görs i par, och hör till kapitel~\ref{c3:ch}. Teorin finns i \secref{c2:app:b} och,
för analys och felsökning, i \secref{c3:app:a}.

\section{Mål}
\label{lab1:sec:goals}
Efter laborationen ska du kunna:
\begin{itemize}
  \item avgöra vilka anslutningar på en tryckknapp som är slutande och vilka som är brytande;
  \item koppla upp ett kontaktnät efter ett strömvägsschema;
  \item realisera AND, OR, NOT, NAND, NOR, XOR och XNOR med kontakter;
  \item syntetisera ett kontaktnät ur en sanningstabell, via ett förenklat uttryck;
  \item analysera ett kontaktnät och ersätta det med ett mindre, likvärdigt nät;
  \item kontrollera ett nät genom att jämföra förutsagda och uppmätta sanningstabeller.
\end{itemize}

\section{Utrustning}
\label{lab1:sec:equipment}
Varje station har:
\begin{itemize}
  \item ett \textbf{DC-aggregat} inställt på 24~V;
  \item två \textbf{tryckknappar}, S1 och S2, med både slutande (NO) och brytande (NC)
    kontaktelement;
  \item en \textbf{glödlampa} för 24~V, H1;
  \item \textbf{kopplingssladdar};
  \item en \textbf{multimeter}, för felsökning.
\end{itemize}

Varje kontaktelement har två anslutningar, märkta med ett nummer. Enligt standarden slutar numret
på \textbf{3 och 4 för ett slutande element} och på \textbf{1 och 2 för ett brytande}, medan den
första siffran anger elementets plats i knappen. Figur~\ref{lab1:fig:station} visar ett exempel;
märkningen kan skilja sig mellan stationerna, och \hyperref[lab1:u1]{uppgift~1} går ut på att ta
reda på hur just din station ser ut.

\bookfigure{lab1_station}{Stationens delar: tryckknapparna S1 och S2, lampan H1 och
aggregatet.}{lab1:fig:station}

Alla uppgifter i laborationen går att göra med \textbf{ett slutande och ett brytande element per
knapp}. Har din station fler, använd dem du behöver.

\section{Säkerhet}
\label{lab1:sec:safety}
\begin{warning}
\begin{itemize}
  \item \textbf{Koppla alltid med aggregatet frånslaget}, och slå på det först när kopplingen är
    kontrollerad mot schemat. 24~V är inte farligt att röra, men en kortslutning ger stora
    strömmar, varma sladdar och ett aggregat som slår ifrån.
  \item \textbf{Varje strömväg ska gå genom lampan.} Koppla aldrig en kontakt direkt mellan
    +24~V och 0~V: när kontakten sluts blir det en kortslutning.
  \item Glödlampan blir varm när den har lyst en stund. Rör inte glaset.
  \item Se också de gemensamma säkerhetsreglerna för laborationerna i förordet.
\end{itemize}
\end{warning}

\section{Förberedelser}
\label{lab1:sec:prep}
Förberedelseuppgifterna ska vara gjorda \textbf{före} passet. Läraren tittar på dem innan ni börjar
koppla, och utan dem hinner man inte klart.

\labprep{F1}{lab1:f1}Skriv sanningstabellen för H1 med ingångarna S1 och S2 för var och en av
funktionerna AND, OR, NAND, NOR, XOR och XNOR, och tabellen för NOT med S1 som enda ingång. En
nedtryckt knapp är \code{1} och en tänd lampa är \code{1}.

\labprep{F2}{lab1:f2}Rita ett strömvägsschema mellan +24~V och 0~V för var och en av funktionerna
i F1, med tryckknappar och lampan H1. Ange för varje kontakt om den är slutande eller brytande.
Lämna plats att fylla i anslutningarnas nummer när du har gjort \hyperref[lab1:u1]{uppgift~1}.

\labprep{F3}{lab1:f3}Läs \hyperref[lab1:u7]{uppgift~7}. Ta fram uttrycket för H1 på SP-form ur
sanningstabellen, förenkla det, och rita det minsta kontaktnätet.

\labprep{F4}{lab1:f4}Läs \hyperref[lab1:u8]{uppgift~8}. Skriv uttrycket för nätet i figuren, ta
fram dess sanningstabell, förenkla uttrycket, och rita det minsta likvärdiga nätet.

\section{Arbetsgång}
\label{lab1:sec:procedure}
Varje uppgift görs på samma sätt, enligt \secref{c3:a4}:
\begin{enumerate}
  \item Ta fram schemat från förberedelserna, och fyll i anslutningarnas nummer.
  \item Fyll i kolumnen \textbf{förutsagt} i uppgiftens tabell.
  \item Koppla med aggregatet frånslaget. Den ena i paret kopplar, den andra kontrollerar mot
    schemat.
  \item Slå på aggregatet, pröva varje kombination av knapparna, och fyll i kolumnen
    \textbf{uppmätt}.
  \item Stämmer kolumnerna? Om inte: felsök enligt \secref{c3:a5}.
  \item \kindtag[fillaccent]{Redovisa} för läraren: schemat, tabellen och ett prov med knapparna.
\end{enumerate}

Byt roll mellan uppgifterna, så att båda i paret har kopplat.

\filbreak
\section{Uppgifterna}
\label{lab1:sec:tasks}

\labtask{Uppgift 1}{Identifiera kontakterna}\phantomsection\label{lab1:u1}
Innan något kan kopplas måste du veta vilka anslutningar som hör till vilket kontaktelement, och om
elementet är slutande eller brytande. Märkningen berättar det, men den ska kontrolleras, eftersom
en felaktigt kopplad kontakt ger en krets som gör precis tvärtom.

Använd lampan som provare: koppla +24~V till en anslutning a, och lampan mellan en annan
anslutning b och 0~V. Om det som sitter mellan a och b leder, lyser lampan.

\bookfigure{lab1_continuity}{Provkopplingen: lampan lyser om det som sitter mellan a och b
leder.}{lab1:fig:continuity}

\begin{parts}
  \item Koppla in ett kontaktelement i taget mellan a och b. Pröva med knappen släppt och nedtryckt,
    och fyll i tabellen, med en rad per element. Har knapparna fler än två element, lägg till rader:

\begin{filltable}{@{}p{3.2em}YYYY@{}}
  \thead{Knapp} & \thead{Anslutningar} & \thead{Lampan, knappen släppt} &
    \thead{Lampan, knappen nedtryckt} & \thead{Slutande eller brytande?} \\
  \midrule
  S1 & & & & \fillrule
  S1 & & & & \fillrule
  S2 & & & & \fillrule
  S2 & & & & \fillrule
  & & & & \fillrule
  & & & & \\
\end{filltable}

  \item Stämmer resultatet med märkningen, enligt regeln att 3--4 är slutande och 1--2 brytande?

  \item Rita av din station som i figur~\ref{lab1:fig:station}, med rätt nummer, och fyll i numren i
    dina scheman från F2.
\end{parts}

\redovisa{tabellen och ritningen.}

\labtask{Uppgift 2}{AND}\phantomsection\label{lab1:u2}
Bygg AND-funktionen, \code{H1 = S1 · S2}, efter ditt schema från F2.

\begin{filltable}{@{}ccYY@{}}
  \thead{S1} & \thead{S2} & \thead{H1 förutsagt} & \thead{H1 uppmätt} \\
  \midrule
  0 & 0 & & \fillrule
  0 & 1 & & \fillrule
  1 & 0 & & \fillrule
  1 & 1 & & \\
\end{filltable}

Nämn ett ställe där en AND-funktion med två knappar används i verkligheten, och förklara varför.

\redovisa{schemat, tabellen och kopplingen.}

\labtask{Uppgift 3}{OR}\phantomsection\label{lab1:u3}
Bygg OR-funktionen, \code{H1 = S1 + S2}.

\begin{filltable}{@{}ccYY@{}}
  \thead{S1} & \thead{S2} & \thead{H1 förutsagt} & \thead{H1 uppmätt} \\
  \midrule
  0 & 0 & & \fillrule
  0 & 1 & & \fillrule
  1 & 0 & & \fillrule
  1 & 1 & & \\
\end{filltable}

\redovisa{schemat, tabellen och kopplingen.}

\labtask{Uppgift 4}{NOT}\phantomsection\label{lab1:u4}
Bygg NOT-funktionen, \code{H1 = S1'}, med S1 som enda knapp.

\begin{filltable}{@{}cYY@{}}
  \thead{S1} & \thead{H1 förutsagt} & \thead{H1 uppmätt} \\
  \midrule
  0 & & \fillrule
  1 & & \\
\end{filltable}

Förklara varför en nödstoppsknapp alltid kopplas med sitt brytande element.

\redovisa{schemat, tabellen och kopplingen.}

\labtask{Uppgift 5}{NAND och NOR}\phantomsection\label{lab1:u5}
\begin{parts}
  \item Bygg NAND-funktionen, \code{H1 = (S1 · S2)'}. Använd De Morgans lag för att skriva om
    uttrycket så att det bara innehåller enskilda, eventuellt inverterade, variabler, och bygg det
    uttrycket.

\begin{filltable}{@{}ccYY@{}}
  \thead{S1} & \thead{S2} & \thead{H1 förutsagt} & \thead{H1 uppmätt} \\
  \midrule
  0 & 0 & & \fillrule
  0 & 1 & & \fillrule
  1 & 0 & & \fillrule
  1 & 1 & & \\
\end{filltable}

  \item Bygg NOR-funktionen, \code{H1 = (S1 + S2)'}, på samma sätt.

\begin{filltable}{@{}ccYY@{}}
  \thead{S1} & \thead{S2} & \thead{H1 förutsagt} & \thead{H1 uppmätt} \\
  \midrule
  0 & 0 & & \fillrule
  0 & 1 & & \fillrule
  1 & 0 & & \fillrule
  1 & 1 & & \\
\end{filltable}

  \item Jämför dina NAND- och NOR-kopplingar med AND- och OR-kopplingarna i uppgift 2 och 3. Vad har
    ändrats, och hur hänger det ihop med De Morgans lagar?
\end{parts}

\redovisa{båda kopplingarna och svaret på c).}

\labtask{Uppgift 6}{XOR och XNOR}\phantomsection\label{lab1:u6}
\begin{parts}
  \item Bygg XOR-funktionen, \code{H1 = S1 · S2' + S1' · S2}. Varje knapp används med både sitt
    slutande och sitt brytande element.

\begin{filltable}{@{}ccYY@{}}
  \thead{S1} & \thead{S2} & \thead{H1 förutsagt} & \thead{H1 uppmätt} \\
  \midrule
  0 & 0 & & \fillrule
  0 & 1 & & \fillrule
  1 & 0 & & \fillrule
  1 & 1 & & \\
\end{filltable}

  \item Gör om kopplingen till XNOR, \code{H1 = S1 · S2 + S1' · S2'}, genom att \textbf{flytta så få
    sladdar som möjligt}. Vilka flyttade du, och varför räcker det?

\begin{filltable}{@{}ccYY@{}}
  \thead{S1} & \thead{S2} & \thead{H1 förutsagt} & \thead{H1 uppmätt} \\
  \midrule
  0 & 0 & & \fillrule
  0 & 1 & & \fillrule
  1 & 0 & & \fillrule
  1 & 1 & & \\
\end{filltable}

  \item XOR-kopplingen kallas ibland trappkoppling. Förklara varför, och vad som skiljer den från en
    riktig trappbelysning med två vippströmbrytare.
\end{parts}

\redovisa{båda kopplingarna och svaren på b) och c).}

\labtask{Uppgift 7}{Syntes från sanningstabell}\phantomsection\label{lab1:u7}
H1 är driftlampan till en maskin. Den ska lysa när maskinens skyddslucka är stängd, eller när
servicenyckeln är vriden, oavsett luckan. I laborationen motsvaras servicenyckeln av S1
(nedtryckt betyder vriden) och luckans givare av S2, som påverkas, alltså är nedtryckt, när luckan
är \textbf{öppen}.

\begin{narrowtable}{@{}ccc@{}}
  \thead{S1} & \thead{S2} & \thead{H1} \\
  \midrule
  0 & 0 & 1 \\
  0 & 1 & 0 \\
  1 & 0 & 1 \\
  1 & 1 & 1 \\
\end{narrowtable}

\begin{parts}
  \item Ta fram uttrycket för H1 på SP-form ur tabellen, och förenkla det (F3). Vilken räknelag
    använder du?

  \item Bygg det minsta kontaktnätet för det förenklade uttrycket, och kontrollera det mot tabellen.

\begin{filltable}{@{}ccYY@{}}
  \thead{S1} & \thead{S2} & \thead{H1 förutsagt} & \thead{H1 uppmätt} \\
  \midrule
  0 & 0 & & \fillrule
  0 & 1 & & \fillrule
  1 & 0 & & \fillrule
  1 & 1 & & \\
\end{filltable}

  \item Hur många kontakter hade nätet behövt om du hade byggt det oförenklade uttrycket, med en
    strömväg per minterm?
\end{parts}

\redovisa{härledningen, kopplingen och svaret på c).}

\labtask{Uppgift 8}{Analys av ett kontaktnät}\phantomsection\label{lab1:u8}
Nätet i figur~\ref{lab1:fig:task8} har konstruerats av någon annan. Din uppgift är att ta reda på
vad det gör, och om det går att göra enklare.

\bookfigure{lab1_task8}{Nätet i uppgift 8.}{lab1:fig:task8}

\begin{parts}
  \item Skriv uttrycket för H1 (F4).

  \item Förutsäg sanningstabellen, bygg nätet precis som det är ritat, och mät.

\begin{filltable}{@{}ccYY@{}}
  \thead{S1} & \thead{S2} & \thead{H1 förutsagt} & \thead{H1 uppmätt} \\
  \midrule
  0 & 0 & & \fillrule
  0 & 1 & & \fillrule
  1 & 0 & & \fillrule
  1 & 1 & & \\
\end{filltable}

  \item Förenkla uttrycket, och bygg det minsta likvärdiga nätet. Mät igen, och kontrollera att
    tabellen är densamma.

\begin{filltable}{@{}ccYY@{}}
  \thead{S1} & \thead{S2} & \thead{H1 förutsagt} & \thead{H1 uppmätt} \\
  \midrule
  0 & 0 & & \fillrule
  0 & 1 & & \fillrule
  1 & 0 & & \fillrule
  1 & 1 & & \\
\end{filltable}

  \item En av kontakterna i det ursprungliga nätet gick att ta bort. Förklara med ord, utan
    räknelagar, varför den aldrig påverkar lampan fast den påverkar strömmen.
\end{parts}

\redovisa{båda näten, tabellerna och förklaringen i d).}

\labtask[Extra]{Uppgift 9}{Fler funktioner}\phantomsection\label{lab1:u9}
\emph{Frivillig uppgift, för den som blir klar i tid.}

Med två knappar finns det 16 olika funktioner, eftersom varje rad i en tabell med fyra rader kan
vara \code{0} eller \code{1}.

\begin{parts}
  \item Skriv upp alla 16 sanningstabellerna, och ange för var och en ett så enkelt uttryck som
    möjligt. Två av dem behöver inga kontakter alls.

  \item Välj fyra funktioner som du inte redan har byggt, bygg dem, och kontrollera dem mot
    tabellerna.

  \item Vilka av de 16 funktionerna behöver både det slutande och det brytande elementet i båda
    knapparna? Vilka klarar sig med ett element per knapp?
\end{parts}

\filbreak
\section{Redovisning}
\label{lab1:sec:report}
Läraren prickar av varje uppgift när den är redovisad. För godkänd laboration krävs uppgift 1--8;
uppgift 9 är frivillig. Båda i paret ska kunna förklara varje koppling.

\begin{filltable}{@{}p{4.6em}Lp{8em}@{}}
  \thead{Uppgift} & \thead{Innehåll} & \thead{Redovisad} \\
  \midrule
  F1--F4 & Förberedelser & \fillrule
  1 & Identifiera kontakterna & \fillrule
  2 & AND & \fillrule
  3 & OR & \fillrule
  4 & NOT & \fillrule
  5 & NAND och NOR & \fillrule
  6 & XOR och XNOR & \fillrule
  7 & Syntes från sanningstabell & \fillrule
  8 & Analys av ett kontaktnät & \fillrule
  9 & Extra (frivillig) & \\
\end{filltable}

Återställ stationen när ni är klara: aggregatet av, sladdarna tillbaka på sin plats.
